% Generated draft: requires mathematical and source review.
\documentclass[11pt]{article}
\usepackage[margin=25mm]{geometry}
\usepackage{fontspec}
\setmainfont{Latin Modern Roman}
\usepackage{amsmath,amssymb,mathtools,unicode-math}
\setmathfont{Latin Modern Math}
\usepackage{enumitem,needspace,xurl,hyperref}
\setlength{\parindent}{0pt}
\setlength{\parskip}{5pt}
\setlength{\emergencystretch}{3em}
\newcommand{\N}{\mathbb{N}}
\newcommand{\Z}{\mathbb{Z}}
\newcommand{\Q}{\mathbb{Q}}
\newcommand{\R}{\mathbb{R}}
\newcommand{\C}{\mathbb{C}}
\newcommand{\F}{\mathbb{F}}
\newcommand{\A}{\mathbb{A}}
\newcommand{\PP}{\mathbb{P}}
\newcommand{\E}{\mathbb{E}}
\newcommand{\eps}{\varepsilon}
\newcommand{\dd}{\,\mathrm{d}}
\begin{document}
\section*{UnsolvedMath ID: 30006618}
\textbf{Source problem code:} catalogue entry
\subsection{Definitions and mathematical statement}
A connected graph \(X\), with vertex set \(V(X)\) and graph distance \(d_X\), is a median graph if, for every three vertices \(x,y,z\), there is a unique vertex in \(I(x,y)\cap I(y,z)\cap I(z,x)\), where
\[
I(x,y)=\{v\in V(X):d_X(x,v)+d_X(v,y)=d_X(x,y)\}.
\]
A regular tree is a connected acyclic graph whose vertices all have the same degree. Under the conventional Cartesian-product interpretation of a product of graphs, two tuples are adjacent precisely when they agree in every coordinate except one, in which their entries are adjacent.

\textbf{Principal question.} Determine which values the travelling-polytope constant \(\operatorname{TP}(X)\) can take on median graphs. In particular, compute \(\operatorname{TP}(X)\) when \(X\) is a product of regular trees. The cleaned statement additionally asks how this constant behaves under products of median graphs.

\textbf{Unresolved definition and scope.} The supplied record does not define \(\operatorname{TP}\). Consequently, it does not specify a self-contained numerical invariant to compute. It also does not specify the product operation, the number of factors, their allowed degrees, or whether bounded degree and unboundedness are assumptions of the question itself. The Cartesian-product interpretation above is conventional and remains to be confirmed against the source. No particular product formula is asserted here.

\textbf{Supplied background assertions.} For an unbounded median graph \(X\) of bounded degree, the extraction reports
\[
\operatorname{TP}(X)\in[1/2,2/3]\cup\{1\},\qquad
\operatorname{TP}(X)=1\ \Longleftrightarrow\ X\text{ is hyperbolic},
\]
and reports \(\operatorname{TP}(X)=1/2\) if \(X\) has infinite cubical dimension. Bounded degree means a uniform finite bound on vertex degrees; cubical dimension conventionally means the supremum of dimensions of cube graphs occurring isometrically in \(X\). The intended definition should be checked. In particular, the combination of bounded degree and infinite cubical dimension needs source review, since an isometrically embedded \(k\)-cube requires vertices of degree at least \(k\). A further extracted formula involves an undefined object printed as \(E^d\) and displaced tokens \(d\) and \(2d-1\); neither that object nor the formula is reliably specified. These are supplied background claims, not independently verified facts.

\textbf{Additional question interleaved in the extraction.} The original statement also contains a passage headed ``Problem 8. Metric actions on median graphs, Genevois.'' For an action \(G\curvearrowright X\), that passage calls the action proper when every vertex stabiliser is finite, and metrically proper when, for an arbitrary fixed \(o\in V(X)\),
\[
\bigl|\{g\in G:d_X(o,go)\le R\}\bigr|<\infty
\quad\text{for every real }R\ge0.
\]
For a subgroup \(N\le G\), the usual permutational wreath product is
\[
(\mathbb Z/2\mathbb Z)\wr_{G/N}G
=\left(\bigoplus_{G/N}\mathbb Z/2\mathbb Z\right)\rtimes G,
\]
where \(G\) permutes the coset coordinates by left translation. The direct sum is viewed as a graph by joining finitely supported configurations that differ at exactly one coordinate.

The passage reports that a proper action of \(G\) on a median graph \(X\) gives a proper action of this wreath product on
\[
\left(\bigoplus_{G/N}\mathbb Z/2\mathbb Z\right)\times X.
\]
It attributes to Chifan and Ioana the assertion that the wreath product is not a-T-menable whenever \(G\) is a-T-menable but \(G/N\) is not. Here a-T-menability conventionally means admitting a proper affine isometric action on a Hilbert space. Treating \(G/N\) as a quotient group requires \(N\) to be normal, an assumption not explicitly stated in this general passage. The extraction then says that a wreath product with free base group \(F_n\) admits a proper median-graph action but no metrically proper one, invoking a theorem of Niblo and Reeves. The condition on the quotient needed for that conclusion is not repeated in the extracted sentence; this is preserved as a qualified source assertion.

The interleaved question asks: if \(N\trianglelefteq F_n\) and \(F_n/N\) does not admit a metrically proper action on any median graph, can
\[
(\mathbb Z/2\mathbb Z)\wr_{F_n/N}F_n
\]
admit a metrically proper action on some median graph? The passage gives \(F_n/N\simeq\operatorname{BS}(1,2)\) as an example of the premise, without supplying a presentation or proof. The rank \(n\) is unspecified. It also asks for a median-geometric understanding of these examples. Whether this separate action problem belongs to the intended permanent record, or was accidentally interleaved with the travelling-polytope question, remains unresolved.
\subsection{Short English statement}
The main question asks for possible travelling-polytope constants of median graphs, especially products of regular trees, and for their behavior under products. The supplied record omits the definition of the constant and several scope assumptions. It also interleaves a separate wreath-product action question. Both the formulation and the supplied background claims require human mathematical review.
\subsection{Sources}
\begin{itemize}
\item Local registry: \texttt{data/unsolvedmath/problems.json}, supplied record with permanent UnsolvedMath ID \(30006618\), exact source code \texttt{catalogue entry}, title ``Travelling-Polytope Constants for Products of Trees.'' The record belongs to the Oberwolfach Reports Open Problems collection and gives proposed year 2026. No proposer is supplied for the travelling-polytope question.
\item Supplied report citation: \emph{Median Geometry and Applications} (2026), \emph{Oberwolfach Reports}, pp. 487--540, DOI \texttt{10.4171/owr/2026/8}. Supplied link: \url{https://doi.org/10.4171/owr/2026/8}. The record describes these reports as documenting mathematical workshops and their open-problem discussions.
\item Supplied literature source, labelled ``OWR report'': \url{https://ems.press/content/serial-article-files/53603}.
\item Literature assessment dated \texttt{2026-08-22}: the supplied curation labels the problem open and says that the extraction contains unresolved travelling-polytope and wreath-product metric-action questions, with no resolution given in the report. This is supplied dated curation, not an independently verified literature assessment. The background explicitly requests independent verification of time-sensitive status and cited claims.
\item The interleaved metric-action passage names Genevois and attributes background results to Chifan and Ioana and to Niblo and Reeves. No separate bibliographic references for those results are supplied, and their precise hypotheses have not been verified.
\item The registry labels the statement \texttt{corrected\_verified} and says that only formula delimiters, notation, grammar, and typography were repaired, preserving mathematical content and scope. These are supplied editorial labels, not independent mathematical verification.
\end{itemize}
\end{document}
